% 来源：sources/2016-2025-期中-甲-历年真题汇编.pdf 第 7 页
\begin{problem}{10 分}{8cm}
设有下列 $n$ 阶行列式：
\[
D_n=\begin{vmatrix}
4&1&0&\cdots&0&0\\
4&4&1&0&\cdots&0\\
0&4&4&1&0&\cdots\\
\vdots&\vdots&\vdots&\vdots&\vdots&\vdots\\
0&\cdots&0&4&4&1\\
0&0&\cdots&0&4&4
\end{vmatrix},
\]
试证明：$D_n=(1+n)2^{n}$。
\end{problem}

\begin{problem}{10 分}{7cm}
设 $\alpha=(1,0,-1)^{T}$，矩阵 $A=\alpha\alpha^{T}$，又设 $n$ 为一正整数，试求 $|2E-(A^{*})^{n}+A^{n}|$。
\end{problem}

\begin{problem}{15 分}{8.5cm}
设有实方阵 $A=\begin{pmatrix}1&3&9\\2&0&6\\-3&1&-7\end{pmatrix}$，$B=\begin{pmatrix}0&a&b\\2&0&6\\-1&1&0\end{pmatrix}$，其中 $a,b$ 为实常数。已知 $r(A)=r(B)$，且线性方程组 $AX=(b,1,0)^{T}$ 有解，试求 $a,b$ 的值。
\end{problem}

\begin{problem}{20 分}{12cm}
当实数 $a$ 取何值时线性方程组
\[
\begin{cases}
-x_1-4x_2+x_3=1,\\
ax_2-3ax_3=3,\\
x_1+3x_2+(a+1)x_3=0
\end{cases}
\]
无解，有唯一解，有无穷解？有解时请求出所有解。
\end{problem}

\begin{problem}{15 分}{8.5cm}
设 $A=\begin{pmatrix}k&1&1&1\\1&k&1&1\\1&1&k&1\\k&k&k&k\end{pmatrix}$，其中 $k$ 为实常数，试求 $A$ 的秩 $r(A)$。
\end{problem}

\begin{problem}{20 分}{12cm}
设 $A,B$ 为 $n$ 阶方阵，试证明：
\begin{enumerate}[label=(\arabic*),leftmargin=2.2em,itemsep=0.25em,topsep=0.3em]
  \item $\operatorname{tr}(A+B)=\operatorname{tr}A+\operatorname{tr}B$；
  \item $\operatorname{tr}(AB)=\operatorname{tr}(BA)$；
  \item 设 $P$ 是一个 $n$ 阶可逆矩阵，则有 $\operatorname{tr}(P^{-1}AP)=\operatorname{tr}A$；
  \item 若 $A=E_{ij}$，其中 $E_{ij}$ 是第 $i$ 行第 $j$ 列处的元素为 $1$、其余元素全部为零的 $n$ 阶方阵，试求 $\operatorname{tr}A$。
\end{enumerate}
\end{problem}

\begin{problem}{5 分}{6.5cm}
设 $A$ 是一个 $n\ (n\ge 2)$ 阶方阵，$A^{*}$ 是 $A$ 的伴随矩阵。若存在 $n$ 维非零列向量 $\alpha$ 使得 $A\alpha=\theta$，其中 $\theta$ 为 $n$ 维零列向量，且非齐次线性方程组 $A^{*}X=\alpha$ 有解，试证明：$r(A)=n-1$。
\end{problem}

\begin{problem}{5 分}{6.5cm}
设 $A$ 为 $n$ 阶可逆矩阵，$\alpha,\beta$ 为两个 $n$ 维列向量，试证明：$|A+\alpha\beta^{T}|=|A|+\beta^{T}A^{*}\alpha$。
\end{problem}
